\section{Projective geometry as finite algebra}\label{sec:geometry}

\subsection{Conventions and elementary identities}
Vectors are columns, and coordinates are numbered $0,1,2$ in algorithms.
Set
\[
 u\cdot v=u_0v_0+u_1v_1+u_2v_2,
\]
\[
 u\times v=(u_1v_2-u_2v_1,\ u_2v_0-u_0v_2,\ u_0v_1-u_1v_0),
 \qquad \br u v w=u\cdot(v\times w)=\det(u,v,w).
\]
Two nonzero triples represent the same projective point when their cross
product is zero; lines have the same representation convention. Over
$\Q$ this agrees with equality up to nonzero scalar multiplication:
choose a nonzero coordinate of the first triple, and the cross-product
equations give the scalar ratio in every coordinate. The relation and
all incidences are invariant under changing representatives.

A point $p$ lies on a line $\ell$ precisely when $\ell\cdot p=0$.
The line joining distinct points $p,q$ is $p\times q$; the point where
distinct lines $\ell,m$ meet is $\ell\times m$. Orthogonality verifies
both incidences. For uniqueness, if $\ell\cdot p=\ell\cdot q=0$, then
\[
 \ell\times(p\times q)=p(\ell\cdot q)-q(\ell\cdot p)=0.
\]
Any nonzero such $\ell$ therefore represents the same line as $p\times q$.
The same identity proves the dual assertion for intersections.

We use determinant multilinearity, sign change on swapping columns,
zero with a repeated column, multiplicativity, and cofactor expansion.
There is no unbounded linear-algebra theorem hidden here. For the fixed
sizes $3,4,5,6$, define
\begin{equation}\label{eq:detperm}
 \det M=\sum_{\sigma\in S_n}\operatorname{sgn}(\sigma)
                   \prod_{i=1}^n M_{i,\sigma(i)}.
\end{equation}
The sum has at most $720$ terms. Multilinearity follows by distributivity.
Swapping columns pairs permutation terms with opposite signs; with equal
columns these terms cancel. Grouping terms by the selected entry in a
row or column gives cofactor expansion. Expanding $\det(UV)$, terms with
repeated selected indices cancel by the same pairing; the remaining
selections are permutations, giving $\det(UV)=\det U\det V$.
These are finite polynomial calculations over any commutative ring.

Two useful consequences, also checked directly from the coordinates, are
\begin{equation}\label{eq:doublecross}
 (u\times v)\times(r\times s)=\br u v s\,r-\br u v r\,s
\end{equation}
and the adjugate identity
\begin{equation}\label{eq:adj3}
 \br u v w\,o
 =(o\cdot u)(v\times w)+(o\cdot v)(w\times u)
                         +(o\cdot w)(u\times v).
\end{equation}
For the latter, the coefficient of each coordinate of $o$ is a cofactor
expansion; off-diagonal coefficients are determinants with a repeated
row. Thus its verification is the fixed-size argument just given.

\begin{lemma}[Incidence gives determinant zero]\label{lem:inc-det}
If $o$ is a nonzero integer triple and
$o\cdot u=o\cdot v=o\cdot w=0$, then $\br u v w=0$.
\end{lemma}
\begin{proof}
The right side of \eqref{eq:adj3} is zero. Select a nonzero coordinate
$o_i$ by testing the three coordinates. Its left coordinate is
$\br u v w\,o_i=0$. Apply \eqref{eq:cancel}.
\end{proof}

\subsection{An actual common line, not just a determinant test}
Define the total function
\begin{equation}\label{eq:axis}
\axis(u,v,w)=
\begin{cases}
 u\times v,&\nz(u\times v),\\
 u\times w,&\neg\nz(u\times v)\text{ and }\nz(u\times w),\\
 (-u_1,u_0,0),&\text{both crosses zero and }(u_0\ne0\text{ or }u_1\ne0),\\
 (u_2,0,0),&\text{otherwise}.
\end{cases}
\end{equation}
Here $\nz(u)$ abbreviates $u_0\ne0\lor u_1\ne0\lor u_2\ne0$.
Each branch uses only integer arithmetic and a fixed number of equality
tests, so this function is primitive recursive.

\begin{lemma}[Constructive common-kernel lemma]\label{lem:axis}
If $u\ne0$ and $\br u v w=0$, then $h=\axis(u,v,w)$ is nonzero and
$h\cdot u=h\cdot v=h\cdot w=0$.
\end{lemma}
\begin{proof}
In the first branch, orthogonality gives the first two equations and
$h\cdot w=\br u v w=0$. In the second, the analogous remaining equation
is $h\cdot v=-\br u v w=0$. Nonzeroness is the respective branch test.
In the third, $h\cdot u=0$, while
$h\cdot v=(u\times v)_2=0$ and
$h\cdot w=(u\times w)_2=0$; the test on $u_0,u_1$ proves $h\ne0$.
In the last branch, $u_0=u_1=0$ and $u_2\ne0$, because $u\ne0$.
Then $h\cdot u=0$ and
$h\cdot v=u_2v_0=(u\times v)_1=0$, with the same argument for $w$.
The component $u_2$ proves $h\ne0$.
\end{proof}
Applied to points, this constructs a common line; applied to lines, it
constructs a common point. It handles repeated nonzero inputs. Thus a
vanishing determinant is not being used as a definition that conceals the
existence assertion in either geometric theorem.

